\section{Los Language Agents de los últimos tres años: ReAct, planificación, uso de herramientas y la ola del código abierto}

Según la entrevista, en los últimos tres años los Language Agents avanzaron más rápido que en las décadas anteriores. Entre los hitos importantes están ReAct, LLM Planner, Mind2Web, Toolformer y AutoGPT, además de una serie de trabajos sobre web agents, computer use agents y robot planning. La transcripción reconoce Toolformer como «2-former»; por el contexto, se refiere a la línea de uso de herramientas del tipo Toolformer de Meta. Esta nota lo trata según el sentido de la entrevista.

\begin{figure}[H]
\centering
\includegraphics[width=0.92\textwidth]{language-agent-stack.png}
\caption{Pila de un Language Agent: del objetivo del usuario a las herramientas, la retroalimentación del entorno y la memoria.}
\end{figure}

\begin{knowledgebox}{Lectura de la figura: los niveles de la pila de un Language Agent}
En el nivel superior está el objetivo del usuario; debajo, el lenguaje natural, los lenguajes de programación y las instrucciones estructuradas; en medio, la planificación y el razonamiento del LLM; más abajo, las interfaces externas: herramientas, API, GUI, CLI, bases de código, etc.; en la base, la retroalimentación del entorno y la memoria. La conclusión que sustenta es esta: un Language Agent no es un truco de prompt, sino una pila de ingeniería de ciclo cerrado.
\end{knowledgebox}

\subsection{ReAct: un insight sencillo pero de gran alcance}

La esencia de ReAct es alternar reasoning y acting: el modelo primero razona, luego actúa, luego observa y vuelve a razonar. Su Yu (苏煜) subraya que muchos trabajos sobre Agents parecen técnicamente sencillos, pero proponer la abstracción correcta en el momento correcto es muy difícil. El valor de ReAct está en que llevó al LLM de la generación de texto en un solo paso a un ciclo de interacción con el entorno.

\noindent\begin{tabularx}{\textwidth}{p{0.22\textwidth}p{0.34\textwidth}X}
\toprule
Trabajo/línea & Problema que resuelve & Relación con el hilo de este episodio \\
\midrule
ReAct & Hacer que el modelo alterne reasoning y acting & Establece el patrón básico de interacción entre el LLM y el entorno. \\
LLM Planner & Usar el LLM para planificación robótica o embodied planning & Aplica el modelo de lenguaje a planes de acción, no solo a conversar. \\
Mind2Web & Construir un benchmark de web/computer use agents & Convierte las tareas web en una capacidad de Agent evaluable. \\
Toolformer & Hacer que el modelo aprenda cuándo invocar herramientas & Conecta el modelo de lenguaje con el ecosistema de API/herramientas externas. \\
AutoGPT & Proyecto temprano de long-horizon agent de código abierto & Muestra lo que el público imagina de los agentes autónomos y también expone sus problemas de fiabilidad. \\
\bottomrule
\end{tabularx}

\subsection{De la proof of concept a la etapa intensiva en recursos}

Su Yu repasa cómo pasó de semantic parsing a language agents y explica por qué cada vez más investigadores dejan la universidad para emprender. Los primeros trabajos sobre Agents eran más bien proofs of concept: una buena idea podía demostrarse con poco costo. A partir de 2025, las ideas de Agents verdaderamente interesantes requieren cada vez más GPU, API, equipos de ingeniería y capacidad de probar y corregir con rapidez, algo que no encaja del todo con la estructura de recursos de la universidad.

\begin{warningbox}{No confundir «el proyecto de código abierto se volvió popular» con «el problema está resuelto»}
Proyectos como AutoGPT y OpenClaw atraen la atención con rapidez porque muestran lo que es posible. Pero lo posible no equivale a un producto fiable. La gestión del estado, la recuperación de errores, el control de costos y los límites de seguridad en tareas de largo plazo son la parte verdaderamente difícil de los Agents.
\end{warningbox}

\subsection{Resumen de la sección}

El desarrollo de los Language Agents en los últimos tres años es el paso de la generación de texto en un solo paso al ciclo cerrado con el entorno. ReAct, la planificación, el uso de herramientas, los web agents y los agentes autónomos de código abierto impulsaron en conjunto la llegada del OpenClaw Moment.
